\chapter{Fractions: All Four Operations}\label{ch:g8:fractions}

This chapter completes the arithmetic of fractions: addition and
subtraction with \emph{any} denominators, multiplication, and ---
the newcomer --- division, which turns out to be a multiplication in
disguise. Signs from \cref{ch:g8:negprod} come along for the ride.

\section{Adding with any denominators}

\begin{method}[Common denominator]\label{met:g8:fractions:add}
To add or subtract $\frac ab$ and $\frac cd$:
\begin{enumerate}
  \item find a common multiple of $b$ and $d$ (the product
        $b \times d$ always works; a smaller one saves effort);
  \item rewrite both fractions with that denominator;
  \item add or subtract the numerators; simplify.
\end{enumerate}
\end{method}

\begin{example}\label{ex:g8:fractions:add}
Compute $\dfrac{5}{6} + \dfrac{3}{4}$. A common multiple of $6$ and
$4$ is $12$:
\[
\frac{5}{6} = \frac{10}{12},
\qquad
\frac{3}{4} = \frac{9}{12},
\qquad
\frac{5}{6} + \frac{3}{4} = \frac{10 + 9}{12} = \frac{19}{12}.
\]
With signs: $\dfrac{2}{3} - \dfrac{7}{5} = \dfrac{10}{15} -
\dfrac{21}{15} = -\dfrac{11}{15}$.
\end{example}

\section{Multiplying fractions}

\begin{theorem}[Product of fractions]\label{thm:g8:fractions:mult}
\[
\frac{a}{b} \times \frac{c}{d} = \frac{a \times c}{b \times d} .
\]
\end{theorem}

\begin{proof}[Why, on a picture]
Take $\frac34$ of $\frac25$ of a square: cut the square in $5$
vertical strips and keep $2$; cut horizontally in $4$ and keep $3$
rows of what remained. The kept part is a grid of $3 \times 2$ small
cells out of $4 \times 5$: fraction $\frac{6}{20}$.
\end{proof}

\begin{omfigure}
\begin{tikzpicture}[scale=0.68]
  \draw[gray!50, thin] (0,0) grid[xstep=1, ystep=1] (5,4);
  \fill[omDef!18] (0,0) rectangle (2,4);
  \fill[omThm!35] (0,0) rectangle (2,3);
  \draw[thick] (0,0) rectangle (5,4);
  \draw[omDef, thick] (2,0) -- (2,4);
  \draw[omThm, thick] (0,3) -- (5,3);
  \node[below, font=\small, omDef] at (1,0) {$\frac25$ kept};
  \node[right, font=\small, omThm] at (5.05,1.5) {$\frac34$ of it};
\end{tikzpicture}

{\small $\frac34 \times \frac25$: the doubly-shaded region is
$3 \times 2 = 6$ cells out of $4 \times 5 = 20$, i.e.\
$\frac{6}{20} = \frac{3}{10}$ --- numerators multiplied, denominators
multiplied.}
\end{omfigure}

\begin{example}\label{ex:g8:fractions:mult}
Simplify \emph{before} multiplying, crossing out common factors:
\[
\frac{7}{12} \times \frac{9}{14}
= \frac{7 \times 9}{12 \times 14}
= \frac{7 \times 3 \times 3}{3 \times 4 \times 2 \times 7}
= \frac{3}{8}
\]
(the factor $7$ and one factor $3$ cancel between top and bottom).
With signs, the rules of \cref{thm:g8:negprod:rules} apply first:
$\left(-\frac23\right) \times \left(-\frac{5}{4}\right) =
+\frac{10}{12} = \frac56$.
\end{example}

\section{Dividing fractions}

\begin{definition}[Inverse]\label{def:g8:fractions:inverse}
The \emph{inverse}\index{inverse} of a nonzero number $x$ is the
number which multiplied by $x$ gives $1$. The inverse of
$\frac ab$ is $\frac ba$, since
$\frac ab \times \frac ba = \frac{ab}{ab} = 1$; the inverse of an
integer $n$ is $\frac1n$.
\end{definition}

\begin{theorem}[Dividing means multiplying by the inverse]\label{thm:g8:fractions:div}
For $c \neq 0$:
\[
\frac{a}{b} \div \frac{c}{d} = \frac{a}{b} \times \frac{d}{c} .
\]
\end{theorem}

\begin{proof}
By definition, the quotient $q$ of $\frac ab$ by $\frac cd$ is the
number with $q \times \frac cd = \frac ab$. Check the candidate
$q = \frac ab \times \frac dc$:
\[
\left(\frac ab \times \frac dc\right) \times \frac cd
= \frac ab \times \left(\frac dc \times \frac cd\right)
= \frac ab \times 1 = \frac ab . \qedhere
\]
\end{proof}

\begin{example}\label{ex:g8:fractions:div}
\[
\frac{3}{5} \div \frac{7}{10}
= \frac{3}{5} \times \frac{10}{7}
= \frac{3 \times 10}{5 \times 7}
= \frac{30}{35} = \frac{6}{7}.
\]
A sanity check: dividing by $\frac{7}{10}$, a number smaller than
$1$, must give a result \emph{larger} than $\frac35$ --- and
$\frac67 > \frac35$ indeed ($\frac{30}{35} > \frac{21}{35}$).
\end{example}

\begin{example}[Fraction bars within fraction bars]\label{ex:g8:fractions:complex}
A ``double-decker'' is just a division written vertically:
\[
\frac{\ \frac{2}{3}\ }{\ \frac{5}{6}\ }
= \frac{2}{3} \div \frac{5}{6}
= \frac{2}{3} \times \frac{6}{5}
= \frac{12}{15} = \frac{4}{5}.
\]
Locate the \emph{main} bar first (the longest one), then apply
\cref{thm:g8:fractions:div}.
\end{example}

\begin{method}[Mixed computations]\label{met:g8:fractions:mixed}
In an expression mixing fractions and the four operations:
\begin{enumerate}
  \item apply the usual priorities (\cref{ch:g7:priorities});
  \item convert every division into a multiplication by the inverse;
  \item decide signs first (\cref{met:g8:negprod:compute}), then
        multiply or find common denominators;
  \item simplify at the earliest opportunity.
\end{enumerate}
\end{method}

\begin{example}\label{ex:g8:fractions:mixed}
\begin{align*}
\frac12 + \frac{2}{3} \times \frac{9}{4}
&= \frac12 + \frac{2 \times 9}{3 \times 4}
&& \text{(product first!)} \\
&= \frac12 + \frac{18}{12} = \frac12 + \frac32
= \frac{4}{2} = 2 .
\end{align*}
\end{example}

\section{Exercises}

\begin{exercise}[$\star$]\label{exo:g8:fractions:1}
Compute and simplify:
\[
\frac{3}{4} + \frac{2}{5}, \qquad
\frac{5}{6} - \frac{3}{8}, \qquad
\frac{7}{10} + \frac{8}{15} .
\]
\end{exercise}

\begin{exercise}[$\star$]\label{exo:g8:fractions:2}
Compute (signs first):
\[
-\frac{2}{7} + \frac{3}{14}, \qquad
\frac{1}{6} - \frac{5}{4}, \qquad
-\frac{3}{5} - \frac{1}{2} .
\]
\end{exercise}

\begin{exercise}[$\star$]\label{exo:g8:fractions:3}
Compute, simplifying before multiplying:
\[
\frac{5}{8} \times \frac{4}{15}, \qquad
\frac{9}{14} \times \frac{7}{3}, \qquad
\left(-\frac{6}{5}\right) \times \frac{10}{9} .
\]
\end{exercise}

\begin{exercise}[$\star$]\label{exo:g8:fractions:4}
Give the inverse of: $\dfrac{3}{7}$; \quad $5$; \quad
$\dfrac{1}{4}$; \quad $-\dfrac{2}{9}$.
\end{exercise}

\begin{exercise}[$\star$]\label{exo:g8:fractions:5}
Compute:
\[
\frac{4}{9} \div \frac{2}{3}, \qquad
\frac{7}{5} \div 14, \qquad
6 \div \frac{3}{4} .
\]
\end{exercise}

\begin{exercise}[$\star$]\label{exo:g8:fractions:6}
Compute the double-deckers:
\[
\frac{\ \frac{3}{4}\ }{\ \frac{9}{8}\ },
\qquad
\frac{\ \frac{5}{6}\ }{\ 10\ },
\qquad
\frac{\ 2\ }{\ \frac{4}{7}\ } .
\]
\end{exercise}

\begin{exercise}[$\star$]\label{exo:g8:fractions:7}
Compute, respecting priorities:
\[
\frac{1}{3} + \frac{1}{2} \times \frac{4}{5},
\qquad
\left(\frac{1}{3} + \frac{1}{2}\right) \times \frac{4}{5} .
\]
\end{exercise}

\begin{exercise}[$\star\star$]\label{exo:g8:fractions:8}
A tank is $\frac{2}{5}$ full. One adds $\frac{1}{3}$ \emph{of the
tank's capacity}. What fraction of the tank is now full? What
fraction is still empty?
\end{exercise}

\begin{exercise}[$\star\star$]\label{exo:g8:fractions:9}
Three quarters of the students of a class passed a test; among those,
two thirds got more than $14$ out of $20$. What fraction of the class
got more than $14$? If that represents $12$ students, how large is
the class?
\end{exercise}

\begin{exercise}[$\star\star$]\label{exo:g8:fractions:10}
A rope of $\frac{15}{2}$ m must be cut into pieces of $\frac{3}{4}$ m
each. How many pieces does one get? (A division of fractions ---
check that the answer is a whole number.)
\end{exercise}

\begin{exercise}[$\star\star$]\label{exo:g8:fractions:11}
Compute
\[
E = \frac{2}{3} - \frac{5}{6} \div \frac{10}{9},
\qquad
F = \left(-\frac{1}{2}\right)^2 \times \frac{8}{3} .
\]
\end{exercise}

\begin{exercise}[$\star\star\star$]\label{exo:g8:fractions:12}
Simplify the expression
\[
1 + \cfrac{1}{1 + \cfrac{1}{1 + 1}}
\]
(start from the innermost fraction and work outward, one bar at a
time).
\end{exercise}
